\subsection{Cohomotoper}

Let H be a quiver, let G be a groupoid, and let \(\varphi\) and \(\psi\) be morphisms of quivers from H to G.

Denote by \(G_{1}\) the quiver defined in the following way: the vertices of \(G_{1}\) are those of G; the arrows of \(G_{1}\) are the elements of the disjoint union of the sets \(\mathrm{Fl(G)}\) and \(\mathrm{Som(H)}\) ; the source map of \(G_{1}\) coincides with that of G on \(\mathrm{Fl(G)}\) and with \(\mathrm{Som}(\varphi)\) on \(\mathrm{Som(H)}\) ; the target map of \(G_{1}\) coincides with that of G on \(\mathrm{Fl(G)}\) and with \(\mathrm{Som}(\psi)\) on \(\mathrm{Som(H)}\) . Denote by \(\alpha_{1}\) the quiver morphism from G to \(G_{1}\) defined by the identity map of \(\mathrm{Som(G)}\) and by the canonical injection of \(\mathrm{Fl(G)}\) into \(\mathrm{Fl(G_{1})}\) . Denote by \(h_{1}\) the canonical injection \(\mathrm{Som(H)} \to \mathrm{Fl(G_{1})}\) .

Consider the free groupoid \(\mathrm{Grp}(G_{1})\) constructed on \(G_{1}\) (II, p. 174, def. 9) and denote by \(\theta_{1}\) the canonical quiver morphism from \(G_{1}\) to \(\mathrm{Grp}(G_{1})\). Finally, denote by \(\mathrm{Coh}(\varphi,\psi)\) the groupoid obtained from \(\mathrm{Grp}(G_{1})\) by contracting the loops (at the origin of \(x\)).

\[
\alpha_ {1} (x) \alpha_ {1} (y) \alpha_ {1} (x y) ^ {- 1},\tag{1}
\]

for every pair \((x,y)\) of composable arrows of G, as well as loops (at \(\varphi(a)\))

\[
\alpha_ {1} (\varphi (f)) h _ {1} (b) \alpha_ {1} (\psi (f)) ^ {- 1} h _ {1} (a) ^ {- 1}\tag{2}
\]

for \(a\) , \(b\) in \(\operatorname{Som}(\mathrm{H})\) and \(f \in \operatorname{Fl}_{ab}(\mathrm{H})\) . Denote by \(\pi: \operatorname{Grp}(\mathrm{G}_1) \to \operatorname{Coh}(\varphi, \psi)\) the canonical morphism; set \(\alpha = \pi \circ \theta_1 \circ \alpha_1\) and \(h = \operatorname{Fl}(\pi \circ \theta_1) \circ h_1\) .

PROPOSITION 2. --- The groupoid Coh(φ,ψ) is generated by the subquiver whose set of vertices is Som(G) and whose set of arrows is the union of the images of the maps Fl(α) and h.

This subquiver being the image of the quiver \(G_{1}\) by the quiver morphism \(\pi \circ \theta_{1}\) , the proposition follows immediately from the construction of \(\mathrm{Coh}(\varphi, \psi)\) .

PROPOSITION 3. --- The quiver morphism \(\alpha\) is a groupoid morphism from G to \(\mathrm{Coh}(\varphi, \psi)\) and the map \(h\) is a homotopy connecting \(\alpha \circ \varphi\) to \(\alpha \circ \psi\) .

The triple \((\mathrm{Coh}(\varphi,\psi),\alpha,h)\) has the following universal property: if \(G'\) is a groupoid, \(\alpha'\) is a morphism of groupoids from G to \(G'\) and \(h'\colon\mathrm{Som}(\mathrm{H})\to\mathrm{Fl}(\mathrm{G}')\) is a homotopy connecting \(\alpha'\circ\varphi\) to \(\alpha'\circ\psi\) , there exists a unique morphism of groupoids \(\eta\colon\mathrm{Coh}(\varphi,\psi)\to\mathrm{G}'\) such that

\[
\alpha^ {\prime} = \eta \circ \alpha \qquad\text{and}\qquad h ^ {\prime} = \operatorname{Fl} (\eta) \circ h.\tag{3}
\]

Given the definition of the groupoid obtained by contracting arrows, the contraction of the loops (1) implies that \(\alpha\) is a morphism of groupoids, and the contraction of the loops (2) implies that \(h\) is a homotopy connecting \(\alpha \circ \varphi\) to \(\alpha \circ \psi\) .

Let \(G'\) , \(\alpha'\) , \(h'\) be as in the statement. Let \(\eta_{1}\) be the quiver morphism from \(G_{1}\) to \(G'\) such that \(\operatorname{Som}(\eta_{1})\) is equal to \(\operatorname{Som}(\alpha')\) and such that \(\operatorname{Fl}(\eta_{1})\) coincides with \(\operatorname{Fl}(\alpha')\) on \(\operatorname{Fl}(G)\) and with \(h'\) on \(\operatorname{Som}(H)\) . There exists a unique groupoid morphism \(\eta_{2}: \operatorname{Grp}(G_{1}) \to G'\) such that \(\eta_{1} = \eta_{2} \circ \theta_{1}\) . Since \(\alpha'\) is a morphism of groupoids and \(h'\) is a homotopy connecting \(\alpha' \circ \varphi\) to \(\alpha' \circ \psi\) , \(\eta_{2}\) defines by passage to the quotient a morphism of groupoids \(\eta\) from \(\operatorname{Coh}(\varphi, \psi)\) to \(G'\) (II, p. 170, prop. 3). This morphism satisfies relations (3) and is the only one (II, p. 185, prop. 2).

DEFINITION 3. --- The groupoid Coh(φ,ψ) is called the cohomotoper of the pair (φ,ψ). One says that α is the canonical morphism from G to Coh(φ,ψ) and that h is the canonical homotopy relating α∘φ to α∘ψ.

We call the frame of the pair \((\varphi,\psi)\) the quiver whose set of vertices is the set of orbits of G, whose set of arrows is the set of connected components of H, and whose source and target maps are induced by \(\varphi\) and \(\psi\) by passing to quotients.

PROPOSITION 4. --- The map Orb(α): Orb(G) → Orb(Coh(φ,ψ)) is surjective, and its fibers are the connected components of the frame of the pair (φ,ψ).

The morphism \(\alpha\) is the composite of the morphisms \(\alpha_{1}\colon\mathrm{G}\to\mathrm{G}_{1}\) , \(\theta_{1}\colon\mathrm{G}_{1}\to\mathrm{Grp}(\mathrm{G}_{1})\) and \(\pi\colon\mathrm{Grp}(\mathrm{G}_{1})\to\mathrm{Coh}(\varphi,\psi)\) . The map \(\theta_{1}\) induces a bijection from the set of connected components of the quiver \(\mathrm{G}_{1}\) to the set of orbits of \(\mathrm{Grp}(\mathrm{G}_{1})\) and the map \(\mathrm{Orb}(\pi)\colon\mathrm{Orb}(\mathrm{Grp}(\mathrm{G}_{1}))\to\mathrm{Orb}(\mathrm{Coh}(\varphi,\psi))\) is bijective (II, p. 170, remark 1). It therefore suffices to prove that the map from \(\mathrm{Orb}(\mathrm{G})\) to \(\pi_{0}(\mathrm{G}_{1})\) induced by \(\alpha_{1}\) is surjective and that its fibers are the connected components of the frame of the pair \((\varphi,\psi)\) . Surjectivity follows from the fact that the map \(\mathrm{Som}(\alpha_{1})\) is the identity map. The equivalence relation in \(\mathrm{Som}(\mathrm{G})\) given by “\(a\) and \(b\) are in the same connected component of \(\mathrm{G}_{1}\)” is generated by the relations “there exists an arrow of G connecting \(a\) to \(b\)” and “there exists a vertex \(h\) of H such that \(\varphi(h)=a\) and \(\psi(h)=b\)”. This equivalence relation is compatible with the map from \(\mathrm{Som}(\mathrm{G})\) to \(\mathrm{Orb}(\mathrm{G})\), and the relation induced on \(\mathrm{Orb}(\mathrm{G})\) is generated by the relation “there exists an orbit \(\eta\) of H such that \(\mathrm{Orb}(\varphi)(\eta)=\alpha\) and \(\mathrm{Orb}(\psi)(\eta)=\beta\)”. This is therefore the relation “\(\alpha\) and \(\beta\) are in the same connected component of the frame of the pair \((\alpha,\beta)\)”.

PROPOSITION 5. --- Let G' be a groupoid, let η, η' be morphisms of groupoids from Coh(φ,ψ) to G' and let k be a map from Som(G) to Fl(G'). For k to be a homotopy connecting η to η', it is necessary and sufficient that the following two conditions be satisfied:

\begin{enumerate}
\def\labelenumi{(\roman{enumi})}
\item
  The map \(k\) is a homotopy connecting \(\eta \circ \alpha\) to \(\eta' \circ \alpha\) ;
\item
  For every vertex a of H, we have
\end{enumerate}

\[
\eta (h (a)) k (\psi (a)) = k (\varphi (a)) \eta^ {\prime} (h (a)).
\]

By definition, for k to be a homotopy connecting \(\eta\) to \(\eta'\) , it is necessary and sufficient that the following two conditions be satisfied (recall that \(\operatorname{Som}(G) = \operatorname{Som}(\operatorname{Coh}(\varphi, \psi))\) ) :

\begin{enumerate}
\def\labelenumi{\alph{enumi})}
\item
  For every vertex \(x\) of \(\operatorname{Coh}(\varphi, \psi)\) , \(k(x)\) connects \(\eta(x)\) to \(\eta'(x)\) ;
\item
  For every pair \((x,y)\) of vertices of \(\mathrm{Coh}(\varphi,\psi)\) and every arrow \(f\in\mathrm{Fl}_{x,y}(\mathrm{Coh}(\varphi,\psi))\) , we have \(\eta(f)k(y)=k(x)\eta'(f)\) .
\end{enumerate}

By prop. 2, it suffices to verify condition \(b\) ) when \(f\) belongs to the image of \(\mathrm{Fl}(\alpha)\) or to that of \(h\) , so that \(b\) ) is equivalent to the conjunction of the two conditions \(c\) ) and \(d\) ) below:

\begin{enumerate}
\def\labelenumi{\alph{enumi})}
\setcounter{enumi}{2}
\item
  For \(x \in \operatorname{Som}(\mathrm{G})\) , \(y \in \operatorname{Som}(\mathrm{G})\) and \(g \in \operatorname{Fl}_{x,y}(\mathrm{G})\) , we have \(\eta(\alpha(g))k(y) = k(x)\eta'(\alpha(g))\) ;
\item
  For every \(a \in \operatorname{Som}(\mathrm{H})\) , we have \(\eta(h(a))k(\psi(a)) = k(\varphi(a))\eta'(h(a))\) . Condition (i) is equivalent to the conjunction of a) and c), and condition (ii) is condition d), whence the corollary.
\end{enumerate}

